\documentclass[10pt,a4paper]{amsart}
\usepackage{amsmath,amssymb,amsthm,graphicx}
% Retained minimal preamble for the Stacks Project chapter varieties.tex.
% Verbatim excerpt of the chapter preamble: Xy-pic setup (native preamble lines 14-18), the native theorem
% declarations with the mathematical macros (native preamble lines 154-217), and the project own no-op
% definitions of the reference/slogan/history display environments (native wrapper.tex lines 1-3) that the
% counted statement uses for its slogan. Copied mechanically; nothing added, removed or redefined.
% For commutative diagrams we use Xy-pic
\usepackage[all]{xy}

% We use 2cell for 2-commutative diagrams.
\xyoption{2cell}
% Theorem environments.
%
\theoremstyle{plain}
\newtheorem{theorem}[subsection]{Theorem}
\newtheorem{proposition}[subsection]{Proposition}
\newtheorem{lemma}[subsection]{Lemma}

\theoremstyle{definition}
\newtheorem{definition}[subsection]{Definition}
\newtheorem{example}[subsection]{Example}
\newtheorem{exercise}[subsection]{Exercise}
\newtheorem{situation}[subsection]{Situation}

\theoremstyle{remark}
\newtheorem{remark}[subsection]{Remark}
\newtheorem{remarks}[subsection]{Remarks}

\numberwithin{equation}{subsection}

% Macros
%
\def\lim{\mathop{\mathrm{lim}}\nolimits}
\def\colim{\mathop{\mathrm{colim}}\nolimits}
\def\Spec{\mathop{\mathrm{Spec}}}
\def\Hom{\mathop{\mathrm{Hom}}\nolimits}
\def\Ext{\mathop{\mathrm{Ext}}\nolimits}
\def\SheafHom{\mathop{\mathcal{H}\!\mathit{om}}\nolimits}
\def\SheafExt{\mathop{\mathcal{E}\!\mathit{xt}}\nolimits}
\def\Sch{\mathit{Sch}}
\def\Mor{\mathop{\mathrm{Mor}}\nolimits}
\def\Ob{\mathop{\mathrm{Ob}}\nolimits}
\def\Sh{\mathop{\mathit{Sh}}\nolimits}
\def\NL{\mathop{N\!L}\nolimits}
\def\CH{\mathop{\mathrm{CH}}\nolimits}
\def\proetale{{pro\text{-}\acute{e}tale}}
\def\etale{{\acute{e}tale}}
\def\QCoh{\mathit{QCoh}}
\def\Ker{\mathop{\mathrm{Ker}}}
\def\Im{\mathop{\mathrm{Im}}}
\def\Coker{\mathop{\mathrm{Coker}}}
\def\Coim{\mathop{\mathrm{Coim}}}

% Boxtimes
%
\DeclareMathSymbol{\boxtimes}{\mathbin}{AMSa}{"02}

%
% Macros for moduli stacks/spaces
%
\def\QCohstack{\mathcal{QC}\!\mathit{oh}}
\def\Cohstack{\mathcal{C}\!\mathit{oh}}
\def\Spacesstack{\mathcal{S}\!\mathit{paces}}
\def\Quotfunctor{\mathrm{Quot}}
\def\Hilbfunctor{\mathrm{Hilb}}
\def\Curvesstack{\mathcal{C}\!\mathit{urves}}
\def\Polarizedstack{\mathcal{P}\!\mathit{olarized}}
\def\Complexesstack{\mathcal{C}\!\mathit{omplexes}}
% \Pic is the operator that assigns to X its picard group, usage \Pic(X)
% \Picardstack_{X/B} denotes the Picard stack of X over B
% \Picardfunctor_{X/B} denotes the Picard functor of X over B
\def\Pic{\mathop{\mathrm{Pic}}\nolimits}
\def\Picardstack{\mathcal{P}\!\mathit{ic}}
\def\Picardfunctor{\mathrm{Pic}}
\def\Deformationcategory{\mathcal{D}\!\mathit{ef}}
\newenvironment{reference}{}{}
\newenvironment{slogan}{}{}
\newenvironment{history}{}{}
\usepackage{cleveref}
\crefname{originallemma}{original lemma}{original lemmas}
\Crefname{originallemma}{Original lemma}{Original lemmas}
\makeatletter
\newlabel{algebra-lemma-geometrically-irreducible-separable-elements}{{0G33}{source}{}{sourcepointer-1}{}}
\newlabel{algebra-lemma-geometrically-irreducible-separable-elements@cref}{{[originallemma][0][]0G33}{[1][0][]source}{}{}{}}
\makeatother
\crefname{originallemma}{original lemma}{original lemmas}
\Crefname{originallemma}{Original lemma}{Original lemmas}
\makeatletter
\newlabel{algebra-lemma-minimal-prime-reduced-ring}{{00EU}{source}{}{sourcepointer-2}{}}
\newlabel{algebra-lemma-minimal-prime-reduced-ring@cref}{{[originallemma][0][]00EU}{[1][0][]source}{}{}{}}
\makeatother
\crefname{originallemma}{original lemma}{original lemmas}
\Crefname{originallemma}{Original lemma}{Original lemmas}
\makeatletter
\newlabel{algebra-lemma-p-ring-map}{{0BRA}{source}{}{sourcepointer-3}{}}
\newlabel{algebra-lemma-p-ring-map@cref}{{[originallemma][0][]0BRA}{[1][0][]source}{}{}{}}
\makeatother
\crefname{originallemma}{original lemma}{original lemmas}
\Crefname{originallemma}{Original lemma}{Original lemmas}
\makeatletter
\newlabel{algebra-lemma-subalgebra-geometrically-irreducible}{{037N}{source}{}{sourcepointer-4}{}}
\newlabel{algebra-lemma-subalgebra-geometrically-irreducible@cref}{{[originallemma][0][]037N}{[1][0][]source}{}{}{}}
\makeatother
\crefname{originallemma}{original lemma}{original lemmas}
\Crefname{originallemma}{Original lemma}{Original lemmas}
\makeatletter
\newlabel{lemma-geometrically-irreducible-local}{{038G}{source}{}{sourcepointer-5}{}}
\newlabel{lemma-geometrically-irreducible-local@cref}{{[originallemma][0][]038G}{[1][0][]source}{}{}{}}
\makeatother
\title{Geometric irreducibility of an irreducible scheme is equivalent to geometric irreducibility of its function field}
\author{The Stacks Project Authors}
\date{stacks-project:varieties.tex:local-snapshot-20261011; GNU FDL 1.2 or later}
\begin{document}
\maketitle
\noindent\textit{Editorial note.} All delivered mathematical text is the author's own, copied from the pinned Stacks Project chapter file varieties.tex; no statement or proof step was rewritten, reordered or removed.\par\medskip
\noindent\textit{Editorial note.} The proof is complete as written and ends with the local criterion applied to X; no step is referred to the reader.\par\medskip
\noindent\textit{Editorial note.} The counted spans carry 7 navigation stubs and no bibliography citation.\par\medskip
\noindent\textit{Editorial note.} The preamble is the chapter's retained minimal preamble (sources/stacks-varieties/retained-preamble.tex lines 1-77), a verbatim excerpt of the native preamble.\par\medskip
\noindent\textit{Editorial note.} The label lemma-geometrically-irreducible-function-field is defined only in varieties.tex, so no cross-chapter restatement had to be weighed for this card.\par\medskip
\noindent\textit{Allowed standard prerequisites (editorial).}\begin{itemize}
\item Geometric irreducibility versus the separable algebraic closure in the field. Algebra, Lemma algebra-lemma-geometrically-irreducible-separable-elements, invoked at proof lines 1457-1459. Supplies the equivalence of conditions (2) and (3).
\item Local criterion for geometric irreducibility of a scheme. Same chapter, Lemma lemma-geometrically-irreducible-local, invoked at proof lines 1464-1466 and 1481-1483. Supplies the passage between geometric irreducibility of X and geometric irreducibility of the rings of affine opens.
\item Subalgebras of geometrically irreducible algebras, and the local ring versus its residue field with locally nilpotent kernel. Algebra, Lemmas algebra-lemma-subalgebra-geometrically-irreducible, algebra-lemma-minimal-prime-reduced-ring and algebra-lemma-p-ring-map, invoked at proof lines 1467-1473 and 1478-1481. Supplies the two directions of the reduction through the local ring at the generic point.
\end{itemize}\medskip
\ifdefined\hypertarget\hypertarget{sourcepointer-1}{}\fi
\noindent\textit{Source reference.} Original lemma 0G33: sources/stacks-chow/tags.txt, source lines 20078--20078. Pointer outside the counted span, invoked once for context; the load-bearing argument of the counted statement is the author's proof reproduced in full in the proof span.\par\medskip
\ifdefined\hypertarget\hypertarget{sourcepointer-2}{}\fi
\noindent\textit{Source reference.} Original lemma 00EU: sources/stacks-chow/tags.txt, source lines 548--548. Pointer outside the counted span, invoked once for context; the load-bearing argument of the counted statement is the author's proof reproduced in full in the proof span.\par\medskip
\ifdefined\hypertarget\hypertarget{sourcepointer-3}{}\fi
\noindent\textit{Source reference.} Original lemma 0BRA: sources/stacks-chow/tags.txt, source lines 14763--14763. Pointer outside the counted span, invoked once for context; the load-bearing argument of the counted statement is the author's proof reproduced in full in the proof span.\par\medskip
\ifdefined\hypertarget\hypertarget{sourcepointer-4}{}\fi
\noindent\textit{Source reference.} Original lemma 037N: sources/stacks-chow/tags.txt, source lines 4206--4206. Pointer outside the counted span, invoked once for context; the load-bearing argument of the counted statement is the author's proof reproduced in full in the proof span.\par\medskip
\ifdefined\hypertarget\hypertarget{sourcepointer-5}{}\fi
\noindent\textit{Source reference.} Original lemma 038G: sources/stacks-varieties/pinned-original-source.tex, source lines 1405--1405. Pointer outside the counted span, invoked once for context; the load-bearing argument of the counted statement is the author's proof reproduced in full in the proof span.\par\medskip
\begin{lemma}
\label{lemma-geometrically-irreducible-function-field}
Let $X$ be an irreducible scheme over the field $k$. Let $\xi \in X$
be its generic point. The following are equivalent:
\begin{enumerate}
\item $X$ is geometrically irreducible over $k$;
\item $\kappa(\xi)$ is geometrically irreducible over $k$;
\item the separable algebraic closure of $k$ in $\kappa(\xi)$ is equal to $k$.
\end{enumerate}
\end{lemma}

\begin{proof}
By Algebra,
Lemma \ref{algebra-lemma-geometrically-irreducible-separable-elements},
(2) and (3) are equivalent.

\medskip\noindent
Assume (1). Recall that $\mathcal{O}_{X, \xi}$ is the filtered
colimit of $\mathcal{O}_X(U)$ where $U$ runs over the nonempty
open affine subschemes of $X$. Combining
Lemma \ref{lemma-geometrically-irreducible-local}
and
Algebra, Lemma \ref{algebra-lemma-subalgebra-geometrically-irreducible}
we see that $\mathcal{O}_{X, \xi}$ is geometrically irreducible over $k$.
Since $\mathcal{O}_{X, \xi} \to \kappa(\xi)$ is a surjection with
locally nilpotent kernel (see
Algebra, Lemma \ref{algebra-lemma-minimal-prime-reduced-ring})
it follows that $\kappa(\xi)$ is geometrically irreducible, see
Algebra, Lemma \ref{algebra-lemma-p-ring-map}.

\medskip\noindent
Assume (2). We may assume that $X$ is reduced. Let $U \subset X$ be
a nonempty affine open. Then $U = \Spec(A)$ where $A$ is a domain
with fraction field $\kappa(\xi)$. Thus $A$ is a $k$-subalgebra
of a geometrically irreducible $k$-algebra. Hence by
Algebra, Lemma \ref{algebra-lemma-subalgebra-geometrically-irreducible}
we see that $A$ is geometrically irreducible over $k$.
By Lemma \ref{lemma-geometrically-irreducible-local} we conclude
that $X$ is geometrically irreducible over $k$.
\end{proof}

\end{document}
